% !TeX root = ../../Book2.tex
% 纲要标题：### 2.3 子模与商模

\section{子模与商模}
\label{b2:sec:0203}

商模的运算已经在\cref{b2:prop:quotient-module}中建立。现在转而考察怎样指定子模，以及模掉一个子模究竟加入了哪些关系。与此相连的另一个问题是：哪些标量在某个元素上，或在整个模上作用为零？这两种零化条件不同，必须分别记录。

% 纲要要点（供目录核对）：
% * 子模生成与关系
% * 零化子与忠实模
% * 模上的零因子和挠

\subsection{子模生成与关系}
\label{b2:subsec:0203:01}

\begin{proposition}\label{b2:prop:generated-submodule}
任意一族 \(M\) 的子模的交仍为子模，空族的交约定为 \(M\)。对任意子集 \(S\subseteq M\)，存在包含 \(S\) 的最小子模，称为 \(S\) \term[shengcheng-zimo]{生成的子模}[submodule generated by a subset]，记为 \(\langle S\rangle_R\)。具体地，
\[
\langle S\rangle_R=
\left\{\,\sum_{j=1}^n r_js_j\ \middle|\ n\ge0,\ r_j\in R,\ s_j\in S\,\right\},
\]
其中空和为零。当 \(S=\{m\}\) 时记为 \(Rm\)；有 \(M=Rm\) 的模称为\term[xunhuan-mo]{循环模}[cyclic module]。
\end{proposition}
\begin{proof}
属于每个子模的元素，其差与任意标量倍仍属于每个子模，所以交为子模。所有包含 \(S\) 的子模形成非空族，因为 \(M\) 在其中；取交便得到最小子模。

右侧的有限和集合包含零及 \(S\)。两个有限和相减仍是有限和，左乘标量 \(a\) 后各系数变为 \(ar_j\)，也仍有同样形式，因此该集合是子模。反过来，任何包含 \(S\) 的子模都包含所有这些有限线性组合，所以它正是最小子模。
\end{proof}
\symidx{\(\langle S\rangle_R\)、\(Rm\)：集合或单个元素生成的子模}

对一族子模 \((L_i)_{i\in I}\)，其\term[zimo-de-he]{和}[sum of submodules]是并集生成的子模，记为 \(\sum_{i\in I}L_i\)。它的每个元素都是有限多个 \(L_i\) 中的元素之和；允许指标集无限，不等于允许在模中作无限加法。例如在 \(\mathbb Z\) 模中，\(4\mathbb Z+6\mathbb Z=2\mathbb Z\)，因为所有和都是偶数，而 \(2=6-4\)。一般并集未必是子模：\(4\) 与 \(6\) 属于 \(4\mathbb Z\cup6\mathbb Z\)，它们的和 \(10\) 却不属于其中。

\begin{definition}\label{b2:def:finitely-generated-module}
若 \(M=Rm_1+\cdots+Rm_n\)，其中 \(n\) 有限，则称 \(M\) 为\term[youxian-shengcheng-mo]{有限生成模}[finitely generated module]，\(m_1,\ldots,m_n\) 为一组生成元。本书也使用简称\term[youxian-mo]{有限模}[finite module]；它不表示底集合只有有限个元素。零模允许由空集生成。
\end{definition}

\(\mathbb Z\) 作为自身上的模由 \(1\) 生成，底集合却无限；\(\mathbb Q\) 作为 \(\mathbb Q\) 模由 \(1\) 生成，作为 \(\mathbb Z\) 模却不有限生成。后一断言可以直接验证：若有限多个有理数的分母都整除正整数 \(d\)，它们的整数线性组合都在 \(d^{-1}\mathbb Z\) 中，而 \(1/(2d)\) 不在其中。因此是否有限生成总是相对于指定标量环而言。

\ProseParagraphBreak
给出生成元并没有给出唯一坐标。称一个等式 \(r_1m_1+\cdots+r_nm_n=0\) 为这组生成元之间的一个关系。若 \(L\le M\)，则在 \(M/L\) 中一个线性组合为零，恰好表示它在 \(M\) 中属于 \(L\)。例如，在
\[
M=\mathbb Z^2/\mathbb Z(2,3)
\]
中，标准向量的像 \(\bar e_1,\bar e_2\) 满足 \(2\bar e_1+3\bar e_2=0\)。这里 \(\mathbb Z^2\) 采用逐坐标运算。映射 \((a,b)\mapsto3a-2b\) 满射到 \(\mathbb Z\)，因为 \((1,1)\) 映到 \(1\)；它的核由 \(3a=2b\) 给出，即 \(\mathbb Z(2,3)\)。由\cref{b2:thm:module-first-isomorphism}，\(M\cong\mathbb Z\)。商中的两个给定生成元因关系而并不独立；这个具体计算预示了随后用关系研究模的方法。

\begin{proposition}\label{b2:prop:cyclic-module-quotient}
设 \(m\in M\)。满同态 \(R\to Rm\)、\(r\mapsto rm\) 的核是左理想
\[
I_m=\{\,r\in R\mid rm=0\,\},
\]
并诱导左 \(R\) 模同构 \(R/I_m\cong Rm\)。所以每个循环左模均是左正则模的一个商。
\end{proposition}
\begin{proof}
映射可加，且把 \(ar\) 送到 \((ar)m=a(rm)\)，故为左模同态，像按定义为 \(Rm\)。核是 \(I_m\)，作为正则模的子模它是左理想。再由\cref{b2:thm:module-first-isomorphism}得到同构。
\end{proof}

这里的 \(R/I_m\) 是商模；只有 \(I_m\) 为双边理想时，陪集乘法才给出商环。通过一个元素描述循环模，不需要把它先做成一个环。

\subsection{零化子与忠实模}
\label{b2:subsec:0203:02}

\begin{definition}\label{b2:def:annihilator-faithful}
对左 \(R\) 模 \(M\) 的子集 \(S\)，定义其\term[linghuazi]{零化子}[annihilator]为
\[
\operatorname{Ann}_R(S)=\{\,r\in R\mid rs=0\text{ 对每个 }s\in S\,\}.
\]
当 \(S=\{m\}\) 时记为 \(\operatorname{Ann}_R(m)\)。若 \(\operatorname{Ann}_R(M)=0\)，称 \(M\) 为\term[zhongshi-mo]{忠实模}[faithful module]。
\end{definition}
\symidx{\(\operatorname{Ann}_R(S)\)：子集 \(S\) 在环 \(R\) 中的零化子}

\cref{b2:def:vector-annihilator}中的 \(U^\circ\) 位于对偶空间，收集在 \(U\) 上为零的线性泛函；这里的 \(\operatorname{Ann}_R(S)\) 位于标量环，收集在 \(S\) 上作用为零的标量。名称相同，但所收集的对象及所在空间不同。

\begin{proposition}\label{b2:prop:annihilator-ideals}
对任意子集 \(S\subseteq M\)，\(\operatorname{Ann}_R(S)\) 是左理想。如果 \(S\) 本身为子模，则它是双边理想。特别地，\(\operatorname{Ann}_R(M)\) 是表示 \(\rho:R\to\operatorname{End}_{\mathbb Z}(M)\) 的核，因而 \(M\) 忠实当且仅当 \(\rho\) 单射。

若 \(I\) 是 \(R\) 的双边理想，原有 \(R\) 作用能通过商环 \(R/I\) 定义为
\[
(r+I)m=rm
\]
当且仅当 \(I\subseteq\operatorname{Ann}_R(M)\)。成立时这个 \(R/I\) 模结构唯一。
\end{proposition}
\begin{proof}
若 \(a,b\) 零化每个 \(s\in S\)，则 \((a-b)s=0\)；任意 \(r\in R\) 满足 \((ra)s=r(as)=0\)，故零化子为左理想。若 \(S\) 是子模，则 \(rs\in S\)，因而 \((ar)s=a(rs)=0\)，又得到右乘封闭性。由\cref{b2:prop:module-representation}的表示公式，\(\rho(a)=0\) 恰好表示 \(aM=0\)，故其核等于全模零化子。

若商环作用有上述形式，取 \(a\in I\)，便有 \(am=(a+I)m=0\)。反过来，若 \(IM=0\)，当 \(r-r'\in I\) 时 \(rm-r'm=0\)，所以作用不依赖代表元。商环的加法、乘法和单位元运算使四条模公理分别从原来的作用下降。由于每个商环标量都是某个 \(r+I\)，规定相容性后没有别的选择。
\end{proof}

每个模都可因此看成 \(R/\operatorname{Ann}_R(M)\) 上的忠实模：若一个陪集零化所有元素，其代表元已属于全模零化子。这个商环只除去对整个模完全不起作用的标量。正则模 \({}_RR\) 忠实，因为零化 \(1_R\) 的标量只能为零；零模的零化子是 \(R\)，所以零模只在 \(R\) 为零环时忠实。

\ProseParagraphBreak
单个元素的零化子可能更大，且未必是双边理想。令 \(R=M_2(k)\) 作用在列向量 \(M=k^2\) 上。一个矩阵零化所有列向量，意味着它的两列都是零，故 \(M\) 忠实。但对 \(e_1=(1,0)^{\mathsf T}\)，
\[
\operatorname{Ann}_R(e_1)
=\left\{\,\begin{pmatrix}0&a\\0&b\end{pmatrix}\ \middle|\ a,b\in k\,\right\}.
\]
它包含 \(E_{12}\)，却不包含 \(E_{12}E_{21}=E_{11}\)，所以不是右理想。忠实性说的是没有一个非零标量零化全部元素，并不要求每个非零元素都具有零零化子。

\subsection{模上的零因子和挠}
\label{b2:subsec:0203:03}

先设 \(R\) 为交换环。给定模 \(M\)，一个标量 \(r\) 是否可以消去，要看映射 \(M\to M\)、\(m\mapsto rm\) 是否单射；这不仅取决于 \(R\)，也取决于 \(M\)。

\begin{definition}\label{b2:def:module-zero-divisor}
若存在 \(0\ne m\in M\) 使 \(rm=0\)，称 \(r\in R\) 为 \(M\) 上的\term[mo-shang-lingyinzi]{零因子}[zero divisor on a module]。本书在这个集合中包括满足条件的零标量：当 \(M\ne0\) 时，\(0\) 是 \(M\) 上的零因子。不是 \(M\) 上零因子的标量称为 \(M\) 上的\term[mo-shang-zhengze-yuan]{正则元}[regular element on a module]。
\end{definition}

例如，\(2\) 不是环 \(\mathbb Z\) 上的零因子，却是 \(\mathbb Z/6\mathbb Z\) 上的零因子，因为 \(2\bar3=0\)。在这个模上，\(2\) 与 \(3\) 都是零因子，但它们的和 \(5\) 可逆地作用，故模上零因子的集合一般不是理想。零模上没有零因子，因为它没有非零元素。

\begin{definition}\label{b2:def:torsion-module}
设 \(R\) 是整环，即非零交换环且没有非零零因子。左 \(R\) 模 \(M\) 的元素 \(m\) 称为\term[naoyuan]{挠元}[torsion element]，若存在 \(0\ne r\in R\) 使 \(rm=0\)。所有挠元的集合记为 \(t_R(M)\)。若 \(t_R(M)=M\)，称 \(M\) 为\term[naomo]{挠模}[torsion module]；若 \(t_R(M)=0\)，称 \(M\) 为\term[wunao-mo]{无挠模}[torsion-free module]。
\end{definition}
\symidx{\(t_R(M)\)：整环 \(R\) 上模 \(M\) 的挠子模}

\begin{proposition}\label{b2:prop:torsion-submodule}
对整环 \(R\) 上的模 \(M\)，\(t_R(M)\) 是子模，称为挠子模，且 \(M/t_R(M)\) 无挠。每个模同态 \(f:M\to N\) 都把 \(t_R(M)\) 映入 \(t_R(N)\)。
\end{proposition}
\begin{proof}
零元由 \(1_R\ne0\) 零化。若 \(ax=0\)、\(by=0\)，其中 \(a,b\ne0\)，则整环条件保证 \(ab\ne0\)，而交换性给出
\[
ab(x-y)=b(ax)-a(by)=0.
\]
故挠元对差封闭。任意 \(r\in R\) 满足 \(a(rx)=r(ax)=0\)，所以也对标量作用封闭。

设 \(0\ne a\in R\)，且 \(a(m+t_R(M))=0\)。这表示 \(am\in t_R(M)\)，因而某个 \(0\ne b\in R\) 满足 \(bam=0\)。由于 \(ba\ne0\)，\(m\) 本身是挠元，其陪集为零，证明商模无挠。最后，若 \(am=0\)，则 \(af(m)=f(am)=0\)，同一个非零标量也零化 \(f(m)\)。
\end{proof}

整环假设在证明中有明确用途：两个非零零化标量相乘后仍须非零。若在带零因子的环上仍用“被某个非零标量零化”定义挠元，它们未必组成子模。取 \(R=M=\mathbb Z/6\mathbb Z\)，\(\bar2\) 与 \(\bar3\) 分别被非零标量 \(\bar3\) 与 \(\bar2\) 零化，但和 \(\bar5\) 是单位，不被任何非零标量零化。本章的挠模结论因此限于整环；在一般环上还可以采用正则标量定义另一种挠条件，使用时须另作约定。

\ProseParagraphBreak
挠性与忠实性也不能混同。\(\mathbb Q/\mathbb Z\) 作为 \(\mathbb Z\) 模是挠模，因为 \(a/b+\mathbb Z\) 被非零整数 \(b\) 零化；它却是忠实的。事实上，对任意非零整数 \(n\)，取 \(q=1/(|n|+1)\)，则 \(nq\notin\mathbb Z\)，故 \(n\) 不能零化整个模。每个元素都有非零零化标量，不表示可以为所有元素选取同一个非零零化标量。

\ProseParagraphBreak
有限生成时则有一个有用的区别。若整环上的挠模由 \(m_1,\ldots,m_n\) 生成，分别选 \(0\ne a_i\) 使 \(a_im_i=0\)，则非零乘积 \(a_1\cdots a_n\) 零化每个生成元，由交换性也零化它们的所有线性组合。于是非零的有限生成挠模不可能忠实。\(\mathbb Q/\mathbb Z\) 的例子说明，这里不能略去有限生成条件。
